% !TeX root = ../tesis.tex
%%%%%%%%%%%%%%%%%%%%%%%%%%%%%%%%%%%%%%%%%%%%%%%%%%%%%%%%%%%%%%%%%%%
% ANEXO B — Documentación Técnica: VFTModel
%%%%%%%%%%%%%%%%%%%%%%%%%%%%%%%%%%%%%%%%%%%%%%%%%%%%%%%%%%%%%%%%%%%

% =============================================================================
%  B.1 — VALIDACIÓN PYDANTIC
% =============================================================================
\section{Contrato de Validación de Datos (Pydantic)}
\label{anx:vft-validacion}

El Cuadro~\ref{tab:vft-campos-validados} presenta la especificación completa
de los campos verificados en el contrato de datos de VFTModel, implementado
en \texttt{src/api/schemas/schemas.py}.

\begin{table}[H]
    \centering
    \caption[Campos validados en el contrato de datos de VFTModel]{Campos
    validados en el contrato de datos de VFTModel.
    Fuente: \texttt{src/api/schemas/schemas.py}.}
    \label{tab:vft-campos-validados}
    \begin{tabularx}{\textwidth}{|l|l|X|}
        \hline
        \textbf{Campo} & \textbf{Tipo} & \textbf{Descripción} \\
        \hline
        \texttt{sistema}                  & Enum (12 valores) &
            Identifica el sistema de transporte \\
        \texttt{tipo\_entidad}            & Enum (2 valores) &
            \texttt{estacion} o \texttt{ruta} \\
        \texttt{jerarquia\_transporte}    & Enum (4 niveles) &
            Jerarquía operativa del sistema \\
        \texttt{derecho\_de\_via}         & Enum (4 valores) &
            Exclusividad del carril de circulación \\
        \texttt{velocidad\_promedio\_kmh} & Flotante &
            Velocidad comercial en km/h \\
        \texttt{frecuencia\_minutos}      & Flotante &
            Intervalo promedio entre servicios \\
        \texttt{sentido}                  & Entero &
            Dirección de la ruta: ida (1), regreso (0) \\
        \texttt{es\_cetram}               & Booleano &
            Si la estación es un nodo de transferencia \textsc{cetram} \\
        \hline
    \end{tabularx}
    \fuentefigura{Elaboración propia con base en
    \texttt{src/api/schemas/schemas.py} (VFTModel, 2026).}
\end{table}

% =============================================================================
%  B.2 — CATÁLOGO COMPLETO DE SISTEMAS
% =============================================================================
\section{Catálogo Completo de Sistemas de Transporte}
\label{anx:vft-catalogo}

El Cuadro~\ref{tab:vft-sistemas} presenta el catálogo completo de los doce
sistemas reconocidos por VFTModel, incluyendo las claves internas del código
y los parámetros de \termino{fallback} definidos en
\path{src/core/models/impedance.py}.

\begin{table}[H]
    \centering
    \caption[Catálogo de sistemas de transporte en VFTModel]{Catálogo de los
    doce sistemas de transporte reconocidos por VFTModel, con claves de código,
    velocidades comerciales de \termino{fallback} y tipo de derecho de vía.
    Fuente: \texttt{src/core/models/impedance.py}.}
    \label{tab:vft-sistemas}
    \begin{tabular}{|l|l|r|l|}
        \hline
        \textbf{Clave} & \textbf{Sistema} & \textbf{Vel. (km/h)} &
        \textbf{Derecho de vía} \\
        \hline
        \texttt{METRO}       & STC Metro                       & 36.0 & Exclusivo  \\
        \texttt{MB}          & Metrobús                        & 16.3 & Confinado  \\
        \texttt{TL}          & Tren Ligero                     & 22.0 & Exclusivo  \\
        \texttt{TROLE}       & Trolebús                        & 18.0 & Compartido \\
        \texttt{RTP}         & Red de Transporte de Pasajeros  & 14.0 & Mixto      \\
        \texttt{CC}          & Corredor Concesionado           & 11.0 & Mixto      \\
        \texttt{CBB}         & Cablebús                        & 20.0 & Exclusivo  \\
        \texttt{MEXICABLE}   & Mexicable (Estado de México)    & 20.0 & Exclusivo  \\
        \texttt{MEXIBÚS}     & Mexibús (Estado de México)      & 16.3 & Confinado  \\
        \texttt{SUB}         & Tren Suburbano                  & 65.0 & Exclusivo  \\
        \texttt{INTERURBANO} & Tren Interurbano México--Toluca & 70.0 & Exclusivo  \\
        \texttt{PUMABUS}     & Pumabús (\textsc{unam})         & 11.0 & Mixto      \\
        \hline
    \end{tabular}
    \fuentefigura{Elaboración propia con base en \texttt{impedance.py}
    (VFTModel, corrida 2026-04-18).}
\end{table}

% =============================================================================
%  B.3 — LOGS DE CONSTRUCCIÓN Y CÓDIGO FUENTE PYTHON
% =============================================================================
\section{Logs de Construcción del Grafo y Código Fuente Python}
\label{anx:vft-codigo}

\subsection*{Logs de construcción (corrida 2026-04-18)}

\begin{figure}[H]
    \centering
    \fontsize{7.5pt}{9pt}\selectfont
    \begin{minipage}{0.95\textwidth}
        \hrule\vspace{1ex}
        \textbf{Logs de Construcción del Grafo (VFT\_Model):}
        \vspace{1ex}
        \begin{verbatim}
            2026-04-18 20:46:41 | INFO | Fase 1 y 2: Extrayendo nodos y trazos base...
            2026-04-18 20:46:43 | INFO | Fase 2 completada: 65220 aristas creadas.
            2026-04-18 20:46:43 | INFO | Grafo Base: 11115 Nodos y 32344 Segmentos.
            2026-04-18 20:46:43 | WARN | Eliminadas 1562 aristas phantom (> 5 km).
            2026-04-18 20:46:43 | INFO | Fase 3: Integrando sistemas (Tolerancia: 85.0m)...
            2026-04-18 20:47:09 | INFO | Se crearon 20482 aristas de Transbordo Peatonal.
            2026-04-18 20:47:09 | INFO | Impedancia aplicada exitosamente a 25197 segmentos.
            2026-04-18 20:47:09 | INFO | Grafo final: 11115 nodos, 45679 aristas.
        \end{verbatim}
        \hrule
    \end{minipage}
\end{figure}

\subsection*{Implementación de la fórmula de Haversine}

\begin{lstlisting}[style=estiloPython,
    caption={Implementación de la fórmula de Haversine.
             \texttt{src/core/models/impedance.py}.},
    label=lst:haversine]
@staticmethod
def haversine(lon1, lat1, lon2, lat2) -> float:
    R = 6_371_000  # Radio de la Tierra en metros
    phi_1        = math.radians(lat1)
    phi_2        = math.radians(lat2)
    delta_phi    = math.radians(lat2 - lat1)
    delta_lambda = math.radians(lon2 - lon1)
    a = math.sin(delta_phi / 2.0) ** 2 + \
        math.cos(phi_1) * math.cos(phi_2) * \
        math.sin(delta_lambda / 2.0) ** 2
    c = 2 * math.atan2(math.sqrt(a), math.sqrt(1 - a))
    return R * c
\end{lstlisting}

\subsection*{Cálculo del Coeficiente de Fricción Vial}

\begin{lstlisting}[style=estiloPython,
    caption={Cálculo del Coeficiente de Fricción Vial.
             \texttt{src/core/models/impedance.py}.},
    label=lst:cf]
class FrictionCalculator:
    BETA_SATURACION_CDMX = 0.759  # TomTom Traffic Index ZMVM

    @classmethod
    def get_friction(cls, derecho_de_via: str) -> float:
        alpha_map = {
            "exclusivo":  0.0,  # Metro, Suburbano, Cablebus
            "confinado":  0.2,  # Metrobus, Mexicable, Mexibus
            "compartido": 0.5,  # Trolebus
            "mixto":      1.0,  # RTP, CC, Pumabus
        }
        alpha = alpha_map.get(derecho_de_via, 1.0)  # default: peor caso
        return 1.0 + (alpha * cls.BETA_SATURACION_CDMX)
\end{lstlisting}

\subsection*{Inyección de impedancia en el grafo}

\begin{lstlisting}[style=estiloPython,
    caption={Inyección del peso de impedancia en cada arista del grafo dirigido.
             \texttt{src/core/models/impedance.py}.},
    label=lst:impedancia]
for u, v, data in self.G.edges(data=True):
    # Transbordos peatonales ya tienen peso asignado por graph_builder
    if data.get("tipo") == "transfer":
        continue

    # d_ij: distancia real acumulada del trazo (o Haversine como fallback)
    d = data.get("distancia_segmento_m") or \
        self.haversine(u[0], u[1], v[0], v[1])

    # v_ij: velocidad comercial -> metros por minuto
    v_kmh  = data.get("velocidad_promedio_kmh") or \
              self.FALLBACK_VELOCIDAD.get(sistema, 15.0)
    v_mmin = v_kmh * (1000.0 / 60.0)

    # CF: coeficiente de friccion segun derecho de via
    cf = FrictionCalculator.get_friction(
             data.get("derecho_de_via", "mixto"))

    # Ecuacion VFT: peso en minutos
    self.G[u][v]["weight"] = round((d / v_mmin) * cf, 4)
\end{lstlisting}

% =============================================================================
%  B.4 — API REST: ENDPOINTS FASTAPI
% =============================================================================
\section{API REST de VFTModel: Endpoints Disponibles}
\label{anx:vft-api}

\begin{table}[H]
    \centering
    \caption[\api{} REST de VFTModel: endpoints disponibles]{Endpoints de la
    \api{} \textsc{rest} de VFTModel y el indicador o acción que producen.
    Fuente: \texttt{src/api/main.py}.}
    \label{tab:vft-endpoints}
    \begin{tabularx}{\textwidth}{|l|l|X|}
        \hline
        \textbf{Endpoint} & \textbf{Método} & \textbf{Resultado} \\
        \hline
        \texttt{/api/v1/network/build-auto}
            & POST & Construye y almacena el grafo en caché \\
        \texttt{/api/v1/network/spatial-coverage}
            & GET  & Cobertura (\%) por alcaldía o municipio \\
        \texttt{/api/v1/network/topological/capillary-strength}
            & GET  & Ranking de fuerza capilar por nodo \\
        \texttt{/api/v1/network/topological/geo-capillary}
            & GET  & Detección de macro-hubs por proximidad espacial \\
        \texttt{/api/v1/network/topological/detour-factor}
            & GET  & Factor de desviación general o por par O-D \\
        \hline
    \end{tabularx}
    \fuentefigura{Elaboración propia con base en \texttt{src/api/main.py}
    (VFTModel, 2026).}
\end{table}
